\documentclass[10pt, aspectratio=169]{beamer}
\usepackage{../beamer_theme_408}

\title{408考研零基础 --- 操作系统}
\subtitle{3-6 进程同步与互斥}
\author{}
\date{}

\begin{document}


\begin{frame}{本节目标}
    \begin{enumerate}
        \item 理解同步与互斥的概念
        \item 掌握临界区及互斥原则
        \item 掌握信号量的概念及PV操作
        \item 能使用信号量解决简单的同步互斥问题
    \end{enumerate}
\end{frame}

\begin{frame}{同步与互斥}
    \begin{columns}
        \begin{column}{0.48\textwidth}
            \textbf{互斥}
            \begin{itemize}
                \item 多个进程\textbf{竞争}同一资源
                \item 同一时刻只允许一个进程访问
                \item 例：打印机只能被一个进程使用
            \end{itemize}
        \end{column}
        \begin{column}{0.48\textwidth}
            \textbf{同步}
            \begin{itemize}
                \item 多个进程\textbf{协作}完成任务
                \item 按约定的先后顺序执行
                \item 例：A生产数据，B处理数据
            \end{itemize}
        \end{column}
    \end{columns}
\end{frame}

\begin{frame}{临界资源与临界区}
    \textbf{临界资源}：一次只允许一个进程使用的资源（如打印机、共享变量）。

    \textbf{临界区}：访问临界资源的代码段。

    \vspace{0.5em}
    \begin{block}{临界区互斥的四个原则}
        \begin{enumerate}
            \item \textbf{空闲让进} —— 资源空闲时允许进程进入
            \item \textbf{忙则等待} —— 资源被占用时其他进程等待
            \item \textbf{有限等待} —— 等待时间有限，避免饥饿
            \item \textbf{让权等待} —— 等待时释放CPU（非忙等）
        \end{enumerate}
    \end{block}
\end{frame}

\begin{frame}{信号量}
    \textbf{信号量（Semaphore）}：一个整数变量 $S$，用于进程同步与互斥。

    \vspace{0.5em}
    \textbf{PV操作}（由Dijkstra提出）：
    \begin{itemize}
        \item $P(S)$：$S = S - 1$；若 $S < 0$，则进程阻塞
        \item $V(S)$：$S = S + 1$；若 $S \le 0$，则唤醒一个阻塞进程
    \end{itemize}
    \vspace{0.5em}
    \begin{alertblock}{注意}
        P操作和V操作是\textbf{原子操作}（不可中断）。
    \end{alertblock}
\end{frame}

\begin{frame}{用信号量实现进程互斥}
    \begin{itemize}
        \item 设置\textbf{互斥信号量} $mutex$，初值为 $1$
    \end{itemize}
    \vspace{0.5em}
    \begin{block}{互斥模式}
        \begin{tabular}{l}
            $P(mutex)$ \\
            \qquad 临界区代码段 \\
            $V(mutex)$ \\
        \end{tabular}
    \end{block}
    \vspace{0.5em}
    \begin{itemize}
        \item $P(mutex)$：若资源空闲，进入临界区；否则阻塞
        \item $V(mutex)$：退出临界区，唤醒等待进程
    \end{itemize}
\end{frame}

\begin{frame}{用信号量实现进程同步}
    同步关系：进程B必须等待进程A完成某操作后才能继续。

    \vspace{0.5em}
    \begin{itemize}
        \item 设置\textbf{同步信号量} $sync$，初值为 $0$
    \end{itemize}
    \vspace{0.5em}
    \begin{block}{同步模式}
        \begin{tabular}{l}
            进程A（先执行）：\quad $V(sync)$ \\
            进程B（后执行）：\quad $P(sync)$ \\
        \end{tabular}
    \end{block}
    \vspace{0.5em}
    \begin{itemize}
        \item A执行后做$V$操作，信号量变为$1$
        \item B执行前做$P$操作，信号量减为$0$，通过
    \end{itemize}
\end{frame}

\begin{frame}{互斥与同步的信号量对比}
    \begin{tabular}{|c|c|c|}
        \hline
        对比项 & 互斥信号量 & 同步信号量 \\
        \hline
        初始值 & $1$ & $0$ \\
        目的 & 保护临界资源 & 保证执行顺序 \\
        使用方式 & P和V成对出现 & 前V后P \\
        信号量含义 & 资源可用数量 & 事件发生次数 \\
        \hline
    \end{tabular}
\end{frame}

\begin{frame}{PV操作注意事项}
    \begin{itemize}
        \item $P$、$V$操作必须\textbf{成对使用}（互斥中）
        \item 多个$P$操作顺序可能导致\textbf{死锁}
        \item PV操作不能在临界区内进行
        \item 信号量不能用于\textbf{传递大量数据}
    \end{itemize}
    \vspace{0.5em}
    \begin{exampleblock}
        只要在临界区前后分别加$P$、$V$，就可以保证同一时刻只有一个进程进入临界区。
    \end{exampleblock}
\end{frame}

\begin{frame}{本节总结}
    \begin{enumerate}
        \item 同步（协作等待）和互斥（竞争资源）是多进程管理的核心问题
        \item 临界区互斥四原则：空闲让进、忙则等待、有限等待、让权等待
        \item 信号量$S$：$P$减1（可能阻塞），$V$加1（可能唤醒）
        \item 互斥信号量初值1，$P\to$临界区$\to V$；同步信号量初值0，前$V$后$P$
    \end{enumerate}
\end{frame}

\end{document}
